\chapter{Harmonic Functions}

In this chapter, we use several different methods to study harmonic functions.
These include mean value properties, fundamental solutions, maximum principles
and the energy method. Four sections in this chapter are relatively independent of
each other.

\section{Mean Value Properties}

We begin this section with the definition of mean value properties. We assume
that $\Omega$ is a connected domain in $\mathbb{R}^n$.

\begin{definition}\label{hanlin:def-mean-value-properties}\dcref{Definition 1.1}\group{harmonic-functions}\source{han-lin-lecture-notes:page-0007}
For $u \in C(\Omega)$, (i) $u$ satisfies the first mean value property if
\[
u(x) = \frac{1}{\omega_n r^{n-1}} \int_{\partial B_r(x)} u(y)\, d\sigma_y \qquad \text{for any } B_r(x) \subset \Omega;
\]
(ii) $u$ satisfies the second mean value property if
\[
u(x) = \frac{n}{\omega_n r^n} \int_{B_r(x)} u(y)\, dy \qquad \text{for any } B_r(x) \subset \Omega,
\]
where $\omega_n$ denotes the surface area of the unit sphere in $\mathbb{R}^n$.
\end{definition}

\begin{remark}\label{hanlin:remark-mean-value-equivalent}\dcref{Remark 1.2}\group{harmonic-functions}\source{han-lin-lecture-notes:page-0007}
These two definitions are equivalent. In fact, if we write (i) as
\[
u(x)r^{n-1} = \frac{1}{\omega_n} \int_{\partial B_r(x)} u(y)\, d\sigma_y,
\]
we may integrate to get (ii). If we write (ii) as
\[
u(x)r^n = \frac{n}{\omega_n} \int_{B_r(x)} u(y)\, dy,
\]
we may differentiate to get (i).
\end{remark}

\begin{remark}\label{hanlin:remark-mean-value-equivalent-forms}\dcref{Remark 1.3}\group{harmonic-functions}\source{han-lin-lecture-notes:page-0007}
We may write the mean value properties in the following equivalent ways:
(i) $u$ satisfies the first mean value property if
\[
u(x) = \frac{1}{\omega_n} \int_{|y|=1} u(x + ry)\, d\sigma_y \qquad \text{for any } B_r(x) \subset \Omega;
\]
(ii) $u$ satisfies the second mean value property if
\[
u(x) = \frac{n}{\omega_n} \int_{|y| \le 1} u(x + ry)\, dy \qquad \text{for any } B_r(x) \subset \Omega.
\]
\end{remark}

\begin{theorem}\label{hanlin:thm-mean-value-max}\dcref{Theorem 1.4}\group{harmonic-functions}\source{han-lin-lecture-notes:page-0008}\uses{hanlin:def-mean-value-properties}
\textit{If $u \in C(\overline{\Omega})$ satisfies the mean value property in $\Omega$, then $u$ assumes its maximum and minimum only on $\partial\Omega$ unless $u$ is constant.}
\end{theorem}

\begin{proof}
We only prove for the maximum. Set
\[
\Sigma = \{ x \in \Omega;\ u(x) = M \equiv \max_{\Omega} u \} \subset \Omega.
\]
It is obvious that $\Sigma$ is relatively closed. Next we show that $\Sigma$ is open. For any $x_0 \in \Sigma$, take $\overline{B_r}(x_0) \subset \Omega$ for some $r > 0$. By the mean value property, we have
\[
M = u(x_0) = \frac{n}{\omega_n r^n}\int_{B_r(x_0)} u(y)\,dy \leq M \frac{n}{\omega_n r^n}\int_{B_r(x_0)} dy = M.
\]
This implies $u = M$ in $B_r(x_0)$. Hence $\Sigma$ is both closed and open in $\Omega$. Therefore either $\Sigma = \emptyset$ or $\Sigma = \Omega$.
\end{proof}

Now we begin to discuss harmonic functions.

\begin{definition}\label{hanlin:def-harmonic}\dcref{Definition 1.5}\group{harmonic-functions}\source{han-lin-lecture-notes:page-0008}
A function $u \in C^2(\Omega)$ is harmonic if $\Delta u = 0$ in $\Omega$.
\end{definition}

\begin{theorem}\label{hanlin:thm-harmonic-mean-value}\dcref{Theorem 1.6}\group{harmonic-functions}\source{han-lin-lecture-notes:page-0008}\uses{hanlin:def-mean-value-properties,hanlin:def-harmonic}
\textit{Let $u \in C^2(\Omega)$ be harmonic in $\Omega$. Then $u$ satisfies the mean value property in $\Omega$.}
\end{theorem}

\begin{proof}
Take any ball $B_r(x) \subset \Omega$. For any $\rho \in (0,r)$, we apply the divergence theorem in $B_\rho(x)$ and get
\[
\int_{B_\rho(x)} \Delta u(y)\,dy = \int_{\partial B_\rho}\frac{\partial u}{\partial \nu}\,d\sigma = \rho^{n-1}\int_{|\omega|=1}\frac{\partial u}{\partial \rho}(x+\rho\omega)\,d\sigma_\omega
\]
\begin{equation}
= \rho^{n-1}\frac{\partial}{\partial \rho}\int_{|\omega|=1} u(x+\rho\omega)\,d\sigma_\omega.
\end{equation}
Hence for the harmonic function $u$, we have for any $\rho \in (0,r)$
\[
\frac{\partial}{\partial \rho}\int_{|\omega|=1} u(x+\rho\omega)\,d\sigma_\omega = 0.
\]
Integrating from $0$ to $r$, we obtain
\[
\int_{|\omega|=1} u(x+r\omega)\,d\sigma_\omega = \int_{|\omega|=1} u(x)\,d\sigma_\omega = u(x)\omega_n,
\]
or
\[
u(x) = \frac{1}{\omega_n}\int_{|\omega|=1} u(x+r\omega)\,d\sigma_\omega = \frac{1}{\omega_n r^{n-1}}\int_{\partial B_r(x)} u(y)\,d\sigma_y.
\]
This finishes the proof.
\end{proof}

\begin{remark}\label{hanlin:remark-mvp-not-smooth}\dcref{Remark 1.7}\group{harmonic-functions}\source{han-lin-lecture-notes:page-0008}\uses{hanlin:thm-harmonic-mean-value}
For a function $u$ satisfying the mean value property, $u$ is not required to be smooth. However a harmonic function is required to be $C^2$. We now prove these two are equivalent.
\end{remark}

\begin{theorem}\label{hanlin:thm-mean-value-harmonic}\dcref{Theorem 1.8}\group{harmonic-functions}\source{han-lin-lecture-notes:page-0008}\uses{hanlin:def-mean-value-properties,hanlin:thm-harmonic-mean-value}
\textit{If $u \in C(\Omega)$ has the mean value property in $\Omega$, then $u$ is smooth and harmonic in $\Omega$.}
\end{theorem}

\begin{proof}
Choose $\varphi \in C_0^\infty(B_1)$ with $\int_{B_1} \varphi = 1$ and $\varphi(x)=\psi(|x|)$, i.e.
\[
\omega_n \int_0^1 r^{n-1}\psi(r)\,dr = 1.
\]
We define $\varphi_\epsilon(z)=\frac{1}{\epsilon^n}\varphi\!\left(\frac{z}{\epsilon}\right)$ for $\epsilon>0$. Now for any $x \in \Omega$, consider $\epsilon<\operatorname{dist}(x,\partial\Omega)$.
Then we have
\[
\int_\Omega u(y)\varphi_\epsilon(y-x)\,dy
= \int_\Omega u(x+y)\varphi_\epsilon(y)\,dy
= \frac{1}{\epsilon^n}\int_{|y|<\epsilon} u(x+y)\varphi\!\left(\frac{y}{\epsilon}\right)\,dy
\]
\[
= \int_{|y|<1} u(x+\epsilon y)\varphi(y)\,dy
\]
\[
= \int_0^1 r^{n-1}\,dr \int_{\partial B_1} u(x+\epsilon rw)\varphi(rw)\,d\sigma_w
\]
\[
= \int_0^1 \psi(r)r^{n-1}\,dr \int_{|w|=1} u(x+\epsilon rw)\,d\sigma_w
\]
\[
= u(x)\omega_n \int_0^1 \psi(r)r^{n-1}\,dr = u(x),
\]
where in the last equality we used the mean value property. Hence we get
\[
u(x)=(\varphi_\epsilon * u)(x) \qquad \text{for any } x \in \Omega_\epsilon = \{y \in \Omega;\ d(y,\partial\Omega)>\epsilon\}.
\]
Therefore, $u$ is smooth. Moreover, by (1) in the proof of Theorem 1.6 and the mean value property, we have for any $B_r(x)\subset \Omega$
\[
\int_{B_r(x)} \Delta u = r^{n-1}\frac{\partial}{\partial r}\int_{|w|=1} u(x+rw)\,d\sigma_w = r^{n-1}\frac{\partial}{\partial r}(\omega_n u(x)) = 0.
\]
This implies $\Delta u = 0$ in $\Omega$.
\end{proof}

\begin{remark}\label{hanlin:remark-mvp-harmonic-smooth}\dcref{Remark 1.9}\group{harmonic-functions}\source{han-lin-lecture-notes:page-0009}\uses{hanlin:thm-mean-value-max,hanlin:thm-harmonic-mean-value,hanlin:thm-mean-value-harmonic}
By combining Theorems 1.4-1.8, we conclude that harmonic functions are smooth and satisfy the mean value property. Hence harmonic functions satisfy the maximum principle, a consequence of which is the uniqueness of solutions of the following Dirichlet problem in a bounded domain
\[
\Delta u = f \text{ in } \Omega
\]
\[
u=\varphi \text{ on } \partial\Omega,
\]
for $f \in C(\Omega)$ and $\varphi \in C(\partial\Omega)$. In general, the uniqueness does not hold for unbounded domains. Consider the following Dirichlet problem in an unbounded domain $\Omega$
\[
\Delta u = 0 \text{ in } \Omega
\]
\[
u=0 \text{ on } \partial\Omega.
\]
First, we consider the case $\Omega=\{x\in \mathbb{R}^n;\ |x|>1\}$. Then we have a nontrivial solution $u$ given by
\[
u(x)=
\begin{cases}
\log |x| & \text{for } n=2;\\
|x|^{2-n}-1 & \text{for } n\ge 3.
\end{cases}
\]
Note $u \to \infty$ as $|x| \to \infty$ for $n = 2$ and $u$ is bounded for $n \ge 3$. Next, we consider the upper half space $\Omega = \{x \in \mathbb{R}^n;\ x_n > 0\}$. Then $u(x) = x_n$ is a nontrivial solution, which is unbounded.
\end{remark}

\begin{lemma}\label{hanlin:lemma-interior-gradient-1}\dcref{Lemma 1.10}\group{harmonic-functions}\source{han-lin-lecture-notes:page-0010}\uses{hanlin:thm-harmonic-mean-value}
\textit{Suppose $u \in C(\overline{B_R(x_0)})$ is harmonic in $B_R(x_0)$. Then there holds}
\[
|Du(x_0)| \le \frac{n}{R} \max_{\overline{B_R(x_0)}} |u|.
\]
\end{lemma}

\begin{proof}
For simplicity, we assume $u \in C^1(\overline{B_R})$. Since $u$ is smooth, then $\Delta(D_{x_i}u) = 0$, i.e., $D_{x_i}u$ is also harmonic in $B_R$. Hence $D_{x_i}u$ satisfies the mean value property. By the divergence theorem, we have
\[
D_{x_i}u(x_0) = \frac{n}{\omega_n R^n}\int_{B_R(x_0)} D_{x_i}u(y)dy = \frac{n}{\omega_n R^n}\int_{\partial B_R(x_0)} u(y)\,\nu_i\,d\sigma_y,
\]
which implies
\[
|D_{x_i}u(x_0)| \le \frac{n}{\omega_n R^n}\max_{\partial B_R(x_0)} |u| \cdot \omega_n R^{n-1} \le \frac{n}{R}\max_{\overline{B_R(x_0)}} |u|.
\]
This finishes the proof. $\square$
\end{proof}

\begin{lemma}\label{hanlin:lemma-interior-gradient-2}\dcref{Lemma 1.11}\group{harmonic-functions}\source{han-lin-lecture-notes:page-0010}\uses{hanlin:thm-harmonic-mean-value}
\textit{Suppose $u \in C(\overline{B_R(x_0)})$ is a nonnegative harmonic function in $B_R(x_0)$. Then there holds}
\[
|Du(x_0)| \le \frac{n}{R}u(x_0).
\]
\end{lemma}

\begin{proof}
As before, by the divergence theorem and the nonnegativeness of $u$, we have
\[
|D_{x_i}u(x_0)| \le \frac{n}{\omega_n R^n}\int_{\partial B_R(x_0)} u(y)\,d\sigma_y = \frac{n}{R}u(x_0),
\]
where in the last equality we used the mean value property.
$\square$
\end{proof}

\begin{corollary}\label{hanlin:corollary-liouville}\dcref{Corollary 1.12}\group{harmonic-functions}\source{han-lin-lecture-notes:page-0010}\uses{hanlin:lemma-interior-gradient-2}
\textit{A harmonic function in $\mathbb{R}^n$ bounded from above or below is constant.}
\end{corollary}

\begin{proof}
Suppose $u$ is a harmonic function in $\mathbb{R}^n$. We prove that $u$ is a constant if $u \ge 0$. In fact for any $x \in \mathbb{R}^n$, we apply Lemma 1.11 to $u$ in $B_R(x)$ and then let $R \to \infty$. We conclude $Du(x) = 0$ for any $x \in \mathbb{R}^n$. $\square$
\end{proof}

\begin{lemma}\label{hanlin:lemma-higher-derivative-estimate}\dcref{Lemma 1.13}\group{harmonic-functions}\source{han-lin-lecture-notes:page-0010}\uses{hanlin:lemma-interior-gradient-1}
\textit{Suppose $u \in C(\overline{B_R(x_0)})$ is harmonic in $B_R(x_0)$. Then there holds for any multi-index $\alpha$ with $|\alpha| = m$}
\[
|D^\alpha u(x_0)| \le \frac{n^m e^{m-1}m!}{R^m}\max_{\overline{B_R(x_0)}} |u|.
\]
\end{lemma}

\begin{proof}
We will prove by an induction on $m \ge 1$. It holds for $m = 1$ by Lemma 1.10. We assume it holds for $m$ and consider $m + 1$. For $0 < \theta < 1$, define $r = (1 - \theta)R \in (0,R)$. We apply Lemma 1.10 to $u$ in $B_r(x_0)$ and get
\[
|D^{m+1}u(x_0)| \le \frac{n}{r}\max_{\overline{B_r(x_0)}} |D^m u|.
\]
By the induction assumption, we have
\[
\max_{\overline{B_r(x_0)}} |D^m u| \le \frac{n^m \cdot e^{m-1} \cdot m!}{(R-r)^m}\max_{\overline{B_R(x_0)}} |u|.
\]
Hence we obtain
\[
|D^{m+1}u(x_0)| \le \frac{n^{m+1}e^{m-1}m!}{R^{m+1}\theta^m(1-\theta)}\max_{\overline{B_R(x_0)}} |u|
= \frac{n^{m+1}e^{m-1}m!}{R^{m+1}\theta^m(1-\theta)}\max_{\overline{B_R(x_0)}} |u|.
\]
By taking $\theta=\dfrac{m}{m+1}$, we have
\[
\frac{1}{\theta^m(1-\theta)}=\left(1+\frac{1}{m}\right)^m(m+1)<e(m+1).
\]
Hence the result is established for any single derivative. For any multi-index $\alpha=(\alpha_1,\cdots,\alpha_n)$, we note $\alpha_1!\cdots\alpha_n!\le (|\alpha|)!$.
\hfill $\square$
\end{proof}

\begin{theorem}\label{hanlin:thm-harmonic-analytic}\dcref{Theorem 1.14}\group{harmonic-functions}\source{han-lin-lecture-notes:page-0011}\uses{hanlin:lemma-higher-derivative-estimate}
\textit{Harmonic functions are analytic.}
\end{theorem}

\begin{proof}
Suppose $u$ is a harmonic function in $\Omega$. For any fixed $x \in \Omega$, take $B_{2R}(x) \subset \Omega$ and $h \in \mathbb{R}^n$ with $|h|\le R$. We have by the Taylor expansion
\[
u(x+h)=u(x)+\sum_{i=1}^{m-1}\frac{1}{i!}\left[\left(h_1\frac{\partial}{\partial x_1}+\cdots+h_n\frac{\partial}{\partial x_n}\right)^i u\right](x)+R_m(h),
\]
where
\[
R_m(h)=\frac{1}{m!}\left[\left(h_1\frac{\partial}{\partial x_1}+\cdots+h_n\frac{\partial}{\partial x_n}\right)^m u\right](x_1+\theta h_1,\ldots,x_n+\theta h_n),
\]
for some $\theta\in(0,1)$. Note that $x+h\in B_R(x)$ for $|h|<R$. Hence by Lemma 1.13, we obtain
\[
|R_m(h)|\le \frac{1}{m!}|h|^m\cdot n^m\cdot \frac{n^m e^{m-1}m!}{R^m}\max_{\overline{B_{2R}(x)}} |u|
\le \left(\frac{|h|n^2e}{R}\right)^m \max_{\overline{B_{2R}(x)}} |u|.
\]
Then for any $h$ with $|h|n^2e<R/2$, $R_m(h)\to 0$ as $m\to\infty$.
\hfill $\square$
\end{proof}

Next we prove the Harnack inequality.

\begin{theorem}\label{hanlin:thm-harnack}\dcref{Theorem 1.15}\group{harmonic-functions}\source{han-lin-lecture-notes:page-0011}\uses{hanlin:thm-harmonic-mean-value}
\textit{Suppose $u$ is a nonnegative harmonic function in $\Omega$. Then for any compact subset $K$ of $\Omega$}
\[
\frac{1}{C}u(y)\le u(x)\le Cu(y)\quad \text{for any } x,y\in K,
\]
\textit{where $C$ is a positive constant depending only on $\Omega$ and $K$.}
\end{theorem}

\begin{proof}
By the mean value property, we can prove easily that there holds for $B_{4R}(x_0)\subset\Omega$
\[
\frac{1}{c}u(y)\le u(x)\le cu(y)\quad \text{for any } x,y\in B_R(x_0),
\]
where $c$ is a positive constant depending only on $n$.
Now for the given compact subset $K$, take $x_1,\ldots,x_N\in K$ such that $\{B_R(x_i)\}$ covers $K$ with $4R<\operatorname{dist}(K,\partial\Omega)$. Then we can choose $C=c^N$.
\hfill $\square$
\end{proof}

\begin{quote}
We finish this section by proving another characterization of harmonic functions. Suppose $u$ is harmonic in $\Omega$. Then we have by integrating by parts
\[
\int_{\Omega} u\Delta\varphi = 0 \qquad \text{for any $\varphi\in C_0^2(\Omega)$.}
\]
The converse is also true.
\end{quote}

\begin{theorem}\label{hanlin:thm-distributional-harmonic}\dcref{Theorem 1.16}\group{harmonic-functions}\source{han-lin-lecture-notes:page-0012}\uses{hanlin:thm-mean-value-harmonic}
\textit{Suppose $u\in C(\Omega)$ satisfies}
\begin{equation}
\int_{\Omega} u\Delta\varphi = 0 \qquad \text{for any $\varphi\in C_0^2(\Omega)$}.
\tag{1}
\end{equation}
\textit{Then $u$ is harmonic in $\Omega$.}
\end{theorem}

\begin{proof}
We claim for any $B_r(x)\subset \Omega$
\begin{equation}
r\int_{\partial B_r(x)} u(y)\,d\sigma_y = n\int_{B_r(x)} u(y)\,dy.
\tag{2}
\end{equation}
Then we have
\[
\frac{d}{dr}\left(\frac{1}{\omega_n r^{n-1}}\int_{\partial B_r(x)} u(y)\,d\sigma_y\right)
= \frac{n}{\omega_n}\frac{d}{dr}\left(\frac{1}{r^n}\int_{B_r(x)} u(y)\,dy\right)
\]
\[
= \frac{n}{\omega_n}\left\{\frac{n}{r^{n+1}}\int_{B_r(x)} u(y)\,dy + \frac{1}{r^n}\int_{\partial B_r(x)} u(y)\,d\sigma_y\right\}=0.
\]
This implies
\[
\frac{1}{\omega_n r^{n-1}}\int_{\partial B_r(x)} u(y)\,dS_y = \text{const.}
\]
This constant is $u(x)$ if we let $r\to 0$. Hence we have
\[
u(x)=\frac{1}{\omega_n r^{n-1}}\int_{\partial B_r(x)} u(y)\,dS_y
\qquad \text{for any $B_r(x)\subset\Omega$.}
\]
Next we prove (2) for $n\geq 3$. For simplicity, we assume $x=0$. Set
\[
\varphi(y,r)=
\begin{cases}
(|y|^2-r^2)^n & |y|\le r\\
0 & |y|>r,
\end{cases}
\]
and then for $k=2,3,\ldots,n$
\[
\varphi_k(y,r)=(|y|^2-r^2)^{n-k}\bigl(2(n-k+1)|y|^2+n(|y|^2-r^2)\bigr)
\qquad \text{for $|y|\le r$.}
\]
A direct calculation shows $\varphi(\cdot,r)\in C_0^2(\Omega)$ and
\[
\Delta_y\varphi(y,r)=
\begin{cases}
2n\varphi_2(y,r) & |y|\le r\\
0 & |y|>r.
\end{cases}
\]
By (1), we have
\[
\int_{B_r(0)} u(y)\varphi_2(y,r)\,dy = 0.
\]
Now we prove if for some $k=2,\ldots,n-1$,
\begin{equation}
\int_{B_r} u(y)\varphi_k(y,r)\,dy = 0,
\tag{3}
\end{equation}
then
(4)
\begin{equation*}
\int_{B_r} u(y)\varphi_{k+1}(y,r)\,dy = 0.
\end{equation*}
In fact, we differentiate (3) with respect to $r$ and get
\begin{equation*}
\int_{\partial B_r} u(y)\varphi_k(y,r)\,dy + \int_{B_r} u(y)\frac{\partial \varphi_k}{\partial r}(y,r)\,dy = 0.
\end{equation*}
For $2 \le k < n$, $\varphi_k(y,r) = 0$ for $|y| = r$. Then we have
\begin{equation*}
\int_{B_r} u(y)\frac{\partial \varphi_k}{\partial r}(y,r)\,dy = 0.
\end{equation*}
We have (4) by noting
\begin{equation*}
\frac{\partial \varphi_k}{\partial r}(y,r) = (-2r)(n-k+1)\varphi_{k+1}(y,r).
\end{equation*}
Therefore by taking $k = n - 1$ in (4), we conclude
\begin{equation*}
\int_{B_r} u(y)\bigl((n+2)|y|^2 - nr^2\bigr)\,dy = 0.
\end{equation*}
Differentiating with respect to $r$ again, we get (2).
\end{proof}

\section{Fundamental Solutions}

\begin{center}
\textbf{1.2. Fundamental Solutions}
\end{center}

We begin this section by seeking for a harmonic function $u$, i.e., $\Delta u = 0$, in $\mathbb R^n$ which depends only on $r = |x - a|$ for some fixed $a \in \mathbb R^n$. By setting $v(r) = u(x)$, we have
\begin{equation*}
v'' + \frac{n-1}{r}v' = 0,
\end{equation*}
and hence
\begin{equation*}
v(r) =
\begin{cases}
c_1 + c_2 \log r & n = 2,\\
c_3 + c_4 r^{2-n} & n \ge 3,
\end{cases}
\end{equation*}
where $c_i$ are constants for $i = 1,2,3,4$. We are interested in functions with a singularity such that
\begin{equation*}
\int_{\partial B_r} \frac{\partial u}{\partial r}\, d\sigma = 1 \qquad \text{for any } r > 0.
\end{equation*}
Hence we set for any fixed $a \in \mathbb R^n$
\begin{equation*}
\Gamma(a,x) =
\begin{cases}
\frac{1}{2\pi}\log |a - x| & \text{for } n = 2,\\
\frac{1}{\omega_n(2-n)}|a - x|^{2-n} & \text{for } n \ge 3.
\end{cases}
\end{equation*}
In summary, for any fixed $a \in \mathbb R^n$, $\Gamma(a,x)$ is harmonic at $x \ne a$, i.e.,
\begin{equation*}
\Delta_x \Gamma(a,x) = 0 \qquad \text{for any } x \ne a,
\end{equation*}
and has a singularity at $x = a$. Moreover, it satisfies
\begin{equation*}
\int_{\partial B_r(a)} \frac{\partial \Gamma}{\partial n_x}(a,x)\, d\sigma_x = 1 \qquad \text{for any } r > 0.
\end{equation*}
The function $\Gamma$ is called the fundamental solution of the Laplacian operator.
Now we prove the Green's identity.

\begin{theorem}\label{hanlin:thm-green-identity}\dcref{Theorem 1.17}\group{harmonic-functions}\source{han-lin-lecture-notes:page-0014}
Suppose $\Omega$ is a bounded domain in $\mathbb{R}^n$ and that
$u \in C^1(\overline{\Omega}) \cap C^2(\Omega)$. Then for any $a \in \Omega$
\[
u(a) = \int_{\Omega} \Gamma(a,x)\Delta u(x)\,dx
- \int_{\partial \Omega}
\left(\Gamma(a,x)\frac{\partial u}{\partial n_x}(x)
- u(x)\frac{\partial \Gamma}{\partial n_x}(a,x)\right)\,d\sigma_x.
\]
\end{theorem}

\begin{remark}\label{hanlin:remark-green-identity}\dcref{Remark 1.18}\group{harmonic-functions}\source{han-lin-lecture-notes:page-0014}
\leavevmode
\begin{enumerate}
\item[(i)] For any $a \in \Omega$, $\Gamma(a,\cdot)$ is integrable in $\Omega$ although it has a singularity.
\item[(ii)] For $a \notin \overline{\Omega}$, the expression in the right hand side is zero.
\item[(iii)] By letting $u=1$, we have
\[
\int_{\partial \Omega}\frac{\partial \Gamma}{\partial n_x}(a,x)\,d\sigma_x = 1
\qquad \text{for any $a \in \Omega$.}
\]
\end{enumerate}
\end{remark}

\begin{proof}
We apply the Green's formula to $u$ and $\Gamma(a,\cdot)$ in $\Omega\setminus B_r(a)$ for small $r>0$ and get
\[
\int_{\Omega \setminus B_r(a)} (\Gamma \Delta u - u\Delta \Gamma)\,dx
= \int_{\partial\Omega}
\left(\Gamma\frac{\partial u}{\partial n} - u\frac{\partial \Gamma}{\partial n}\right)\,d\sigma_x
- \int_{\partial B_r(a)}
\left(\Gamma\frac{\partial u}{\partial n} - u\frac{\partial \Gamma}{\partial n}\right)\,d\sigma_x.
\]
By Noting $\Delta \Gamma = 0$ in $\Omega \setminus B_r(a)$, we have
\[
\int_{\Omega}\Gamma \Delta u\,dx
= \int_{\partial\Omega}\left(\Gamma\frac{\partial u}{\partial n} - u\frac{\partial \Gamma}{\partial n}\right)\,d\sigma_x
- \lim_{r\to 0}\int_{\partial B_r(a)}
\left(\Gamma\frac{\partial u}{\partial n} - u\frac{\partial \Gamma}{\partial n}\right)\,d\sigma_x.
\]
For $n \geq 3$, we get by the definition of $\Gamma$
\[
\left|\int_{\partial B_r(a)} \Gamma \frac{\partial u}{\partial n}\,d\sigma\right|
= \left|\frac{1}{(2-n)\omega_n}r^{2-n}\int_{\partial B_r(a)}\frac{\partial u}{\partial n}\,d\sigma\right|
\leq \frac{r}{n-2}\sup_{\partial B_r(a)}|Du| \to 0 \qquad \text{as } r\to 0,
\]
and
\[
\int_{\partial B_r(a)}u\frac{\partial \Gamma}{\partial n}\,dS
= \frac{1}{\omega_n r^{n-1}}\int_{\partial B_r(a)}u\,dS \to u(a) \qquad \text{as } r\to 0.
\]
For $n=2$, we get the same conclusion similarly.
\end{proof}

\begin{remark}\label{hanlin:remark-local-green-gradient}\dcref{Remark 1.19}\group{harmonic-functions}\source{han-lin-lecture-notes:page-0014}\uses{hanlin:thm-green-identity}
We may employ the local version of the Green's identity to get gradient estimates without using the mean value property. Suppose $u \in C(\overline{B_1})$ is harmonic in $B_1$. For any fixed $0<r<R<1$, choose a cut-off function $\varphi \in C_0^{\infty}(B_R)$ such that $\varphi = 1$ in $B_r$ and $0 \leq \varphi \leq 1$. Apply the Green's formula to $u$ and $\varphi\Gamma(a,\cdot)$ in $B_1\setminus B_{\rho}(a)$ for $a \in B_r$ and $\rho$ small enough. We proceed as in the proof of Theorem 1.17 to obtain
\[
u(a) = - \int_{r<|x|<R} u(x)\,\Delta_x(\varphi(x)\Gamma(a,x))\,dx
\qquad \text{for any } a \in B_r.
\]
Hence, we get
\[
\sup_{B_{1/2}} |u|
\leq c\left(\int_{B_1}|u|^p\right)^{\frac{1}{p}},
\]
and

\[
\sup_{B_{1/2}} |Du| \leq c \max_{B_1} |u|,
\]

where $c$ is a constant depending only on $n$.
\end{remark}

Now we begin to discuss the Green’s functions. Suppose $\Omega$ is a bounded domain in $\mathbb{R}^n$. For any $u \in C^1(\overline{\Omega}) \cap C^2(\Omega)$, we have by Theorem 1.17 for any $x \in \Omega$

\[
u(x) = \int_{\Omega} \Gamma(x,y)\Delta u(y)\,dy - \int_{\partial \Omega} \left(\Gamma(x,y)\frac{\partial u}{\partial n_y}(y) - u(y)\frac{\partial \Gamma}{\partial n_y}(x,y)\right)\,d\sigma_y.
\]

If $u$ solves the following Dirichlet boundary value problem

\[
(\ast)\qquad
\begin{cases}
\Delta u = f & \text{in } \Omega \\
u = \varphi & \text{on } \partial \Omega
\end{cases}
\]

for some $f \in C(\overline{\Omega})$ and $\varphi \in C(\partial \Omega)$, then $u$ can be expressed in terms of $f$ and $\varphi$, with one unknown term. We intend to eliminate this term by adjusting $\Gamma$.

For any fixed $x \in \Omega$, consider

\[
\gamma(x,y) = \Gamma(x,y) + \Phi(x,y),
\]

for some $\Phi(x,\cdot) \in C^2(\overline{\Omega})$ with $\Delta_y \Phi(x,y) = 0$ in $\Omega$. Then Theorem 1.17 can be expressed as follows for any $x \in \Omega$

\[
u(x) = \int_{\Omega} \gamma(x,y)\Delta u(y)\,dy - \int_{\partial \Omega} \left(\gamma(x,y)\frac{\partial u}{\partial n_y}(y) - u(y)\frac{\partial \gamma}{\partial n_y}(x,y)\right)\,d\sigma_y,
\]

since the extra $\Phi(x,\cdot)$ is harmonic. Now by choosing $\Phi$ appropriately, we are led to the important concept of Green’s functions.

For each fixed $x \in \Omega$, choose $\Phi(x,\cdot) \in C^1(\overline{\Omega}) \cap C^2(\Omega)$ such that

\[
\begin{cases}
\Delta_y \Phi(x,y) = 0 & \text{for } y \in \Omega \\
\Phi(x,y) = -\Gamma(x,y) & \text{for } y \in \partial \Omega.
\end{cases}
\]

Denote by $G(x,y)$ the resulting $\gamma(x,y)$, which is called the Green’s function. If such a $G$ exists, then the solution $u$ of the Dirichlet problem $(\ast)$ can be expressed by

\[
(**)\qquad
u(x) = \int_{\Omega} G(x,y)f(y)\,dy + \int_{\partial \Omega} \varphi(y)\frac{\partial G}{\partial n_y}(x,y)\,d\sigma_y.
\]

Note that the Green’s function $G(x,y)$ is defined as a function of $y \in \overline{\Omega}$ for each fixed $x \in \Omega$. Now we discuss some properties of $G$ as a function of $x$ and $y$. We first note by the maximum principle that the Green’s function is unique, since the difference of two Green’s functions is harmonic in $\Omega$ with a zero boundary value.

\begin{lemma}\label{hanlin:lemma-green-symmetric}\dcref{Lemma 1.20}\group{harmonic-functions}\source{han-lin-lecture-notes:page-0015,han-lin-lecture-notes:page-0016}\uses{hanlin:thm-green-identity}
\emph{The Green’s function $G(x,y)$ is symmetric in $\Omega \times \Omega$, i.e., $G(x,y) = G(y,x)$ for $x \neq y \in \Omega$.}
\end{lemma}

\begin{proof}
For any $x_1, x_2 \in \Omega$ with $x_1 \neq x_2$, take $r > 0$ small such that $B_r(x_1) \cap B_r(x_2) = \emptyset$. Set $G_1(y) = G(x_1,y)$ and $G_2(y) = G(x_2,y)$. We apply the Green’s
formula in $\Omega \setminus B_r(x_1)\cup B_r(x_2)$ and get
\[
\int_{\Omega\setminus B_r(x_1)\cup B_r(x_2)}
(G_1\Delta G_2-G_2\Delta G_1)
=\int_{\partial\Omega}\left(G_1\frac{\partial G_2}{\partial n}-G_2\frac{\partial G_1}{\partial n}\right)\dd\sigma
\]
\[
-\int_{\partial B_r(x_1)}\left(G_1\frac{\partial G_2}{\partial n}-G_2\frac{\partial G_1}{\partial n}\right)\dd\sigma
-\int_{\partial B_r(x_2)}\left(G_1\frac{\partial G_2}{\partial n}-G_2\frac{\partial G_1}{\partial n}\right)\dd\sigma.
\]
Since $G_i$ is harmonic for $y\neq x_i$, $i=1,2$, and vanishes on $\partial\Omega$, we have
\[
\int_{\partial B_r(x_1)}\left(G_1\frac{\partial G_2}{\partial n}-G_2\frac{\partial G_1}{\partial n}\right)\dd\sigma
+\int_{\partial B_r(x_2)}\left(G_1\frac{\partial G_2}{\partial n}-G_2\frac{\partial G_1}{\partial n}\right)\dd\sigma=0.
\]
Note that the left hand side has the same limit as the left hand side in the following
as $r\to 0$
\[
\int_{\partial B_r(x_1)}\left(\Gamma\frac{\partial G_2}{\partial n}-G_2\frac{\partial \Gamma}{\partial n}\right)\dd\sigma
+\int_{\partial B_r(x_2)}\left(G_1\frac{\partial \Gamma}{\partial n}-\Gamma\frac{\partial G_1}{\partial n}\right)\dd\sigma=0.
\]
On the other hand, we have
\[
\int_{\partial B_r(x_1)}\Gamma\frac{\partial G_2}{\partial n}\dd\sigma\to 0,\qquad
\int_{\partial B_r(x_2)}\Gamma\frac{\partial G_1}{\partial n}\dd\sigma\to 0
\qquad \text{as }r\to 0,
\]
and
\[
\int_{\partial B_r(x_1)}G_2\frac{\partial \Gamma}{\partial n}\dd\sigma\to G_2(x_1),\qquad
\int_{\partial B_r(x_2)}G_1\frac{\partial \Gamma}{\partial n}\dd\sigma\to G_1(x_2)
\qquad \text{as }r\to 0.
\]
This implies $G_2(x_1)-G_1(x_2)=0$, or $G(x_2,x_1)=G(x_1,x_2)$.
\hfill $\square$
\end{proof}

\begin{proposition}\label{hanlin:prop-green-sign}\dcref{Proposition 1.21}\group{harmonic-functions}\source{han-lin-lecture-notes:page-0016}\uses{hanlin:thm-mean-value-max}
Let $G$ be the Green's function on $\Omega$. Then for any $x,y\in\Omega$
with $x\neq y$
\[
0>G(x,y)>\Gamma(x,y)\qquad \text{for }n\ge 3,
\]
\[
0>G(x,y)>\Gamma(x,y)-\frac{1}{2\pi}\log \diam(\Omega)\qquad \text{for }n=2.
\]
\end{proposition}
\begin{proof}
Fix an $x\in\Omega$ and write $G(y)=G(x,y)$. Since $\lim_{y\to x}G(y)=-\infty$, there
exists an $r>0$ such that $G(y)<0$ in $B_r(x)$. Note that $G$ is harmonic in $\Omega\setminus B_r(x)$
with $G=0$ on $\partial\Omega$ and $G<0$ on $\partial B_r(x)$. The maximum principle implies $G(y)<0$
in $\Omega\setminus B_r(x)$. Next, we discuss the other part of the inequality. Recall the definition
of the Green's function
\[
G(x,y)=\Gamma(x,y)+\Phi(x,y),
\]
where
\[
\Delta \Phi=0 \qquad \text{in }\Omega
\]
\[
\Phi=-\Gamma \qquad \text{on }\partial\Omega.
\]
For $n\ge 3$, we have
\[
\Gamma(x,y)=\frac{1}{(2-n)\omega_n}|x-y|^{2-n}<0 \qquad \text{for } y\in\partial\Omega,
\]
which implies $\Phi(x,\cdot)>0$ on $\partial\Omega$. By the maximum principle, we have $\Phi>0$ in $\Omega$.
For $n=2$, we have
\[
\Gamma(x,y)=\frac{1}{2\pi}\log |x-y|\le \frac{1}{2\pi}\log \diam(\Omega)\qquad \text{for } y\in\partial\Omega.
\]

Hence the maximum principle implies $\Phi > -\frac{1}{2\pi}\log \diam(\Omega)$ in $\Omega$.
\hfill $\square$
\end{proof}

Now, we calculate Green’s functions explicitly for some special domains.

\begin{theorem}\label{hanlin:thm-green-ball}\dcref{Theorem 1.22}\group{harmonic-functions}\source{han-lin-lecture-notes:page-0017}
\emph{The Green’s function for the ball $B_R \subset \mathbb{R}^n$ is given by}
\emph{(i) for $n \ge 3$}

\[
G(x,y)=\frac{1}{(2-n)\omega_n}\left(|x-y|^{2-n}-\left|\frac{R}{|x|}x-\frac{|x|}{R}y\right|^{2-n}\right);
\]

\emph{(ii) for $n=2$}

\[
G(x,y)=\frac{1}{2\pi}\left(\log |x-y|-\log \left|\frac{R}{|x|}x-\frac{|x|}{R}y\right|\right).
\]
\end{theorem}

\begin{proof}
Fix an $x\neq 0$ with $|x|<R$, and consider $X\in \mathbb{R}^n\setminus \bar{B}_R$ given by
\[
X=\frac{R^2}{|x|^2}x.
\]
In other words, $X$ and $x$ are reflexive of each other with respect to the sphere $\partial B_R$. Note that the map $x\mapsto X$ is conformal, i.e., preserves angles. For $|y|=R$, we have by the similarity of triangles
\[
\frac{|x|}{R}=\frac{R}{|X|}=\frac{|y-x|}{|y-X|},
\]
which implies
\begin{equation}
|y-x|=\frac{|x|}{R}|y-X|=\left|\frac{|x|}{R}y-\frac{R}{|x|}x\right|
\qquad \text{for any } y\in \partial B_R.
\tag{1}
\end{equation}
Therefore, in order to have a zero boundary value, we take for $n\ge 3$
\[
G(x,y)=\frac{1}{(2-n)\omega_n}\left(\frac{1}{|x-y|^{n-2}}-\left(\frac{R}{|x|}\right)^{n-2}\frac{1}{|y-X|^{n-2}}\right).
\]
The case $n=2$ is similar.
\end{proof}

Next, we calculate the normal derivative of the Green’s function on the sphere.

\begin{corollary}\label{hanlin:cor-green-normal-derivative}\dcref{Corollary 1.23}\group{harmonic-functions}\source{han-lin-lecture-notes:page-0017}\uses{hanlin:thm-green-ball}
\emph{Suppose $G$ is the Green’s function in $B_R$. Then}
\[
\frac{\partial G}{\partial n_y}(x,y)=\frac{R^2-|x|^2}{\omega_n R|x-y|^n}
\qquad \text{for any } x\in B_R \text{ and } y\in \partial B_R.
\]
\end{corollary}
\begin{proof}
We only consider the case $n\ge 3$. Recall with $X=R^2x/|x|^2$
\[
G(x,y)=\frac{1}{(2-n)\omega_n}\left(|x-y|^{2-n}-\left(\frac{R}{|x|}\right)^{n-2}|y-X|^{2-n}\right),
\]
for any $x\in B_R$ and $y\in \partial B_R$. Hence we have for such $x$ and $y$
\[
D_{y_i}G(x,y)=-\frac{1}{\omega_n}\left(\frac{x_i-y_i}{|x-y|^n}-\left(\frac{R}{|x|}\right)^{n-2}\frac{X_i-y_i}{|X-y|^n}\right)
=\frac{y_i}{\omega_n R}\frac{R^2-|x|^2}{|x-y|^n},
\]
by (1) in the proof of Theorem 1.22. With $n_i=\frac{y_i}{R}$ for $|y|=R$, we obtain
\[
\frac{\partial G}{\partial n_y}(x,y)=\sum_{i=1}^n n_iD_{y_i}G(x,y)=\frac{1}{\omega_n R}\frac{R^2-|x|^2}{|x-y|^n}.
\]
This finishes the proof.
\end{proof}

\begin{quote}
Denote by $K(x,y)$ the function in Corollary 1.23 for $x\in \Omega, y\in \partial \Omega$. It is called
the Poisson kernel and has the following properties:
\begin{align*}
\text{(K1)}\quad &K(x,y) \text{ is smooth for } x\neq y.\\
\text{(K2)}\quad &K(x,y)>0 \text{ for } |x|<R \text{ and } |y|=R.\\
\text{(K3)}\quad &\text{For any fixed } |x_0|=R,\ \lim_{x\to x_0,\,|x|<R}K(x,y)=0 \text{ uniformly in } y \text{ for }\\
&|y-x_0|>\delta>0.\\
\text{(K4)}\quad &\Delta_x K(x,y)=0 \text{ for } |x|<R \text{ and } |y|=R.\\
\text{(K5)}\quad &\int_{|y|=R}K(x,y)dS_y=1 \text{ for any } |x|<R.
\end{align*}
Here (K1), (K2) and (K3) follow easily from the explicit expression for $K$ in
Corollary 1.23 and (K4) follows easily from the definition $K(x,y)=\partial_{n_y}G(x,y)$ and
the fact that $G(x,y)$ is harmonic in $x$. An easy derivation of (K5) is based on (**).
By taking a $C^2(\overline{B_R})$ harmonic function $u$ in (**) , we conclude
\[
u(x)=\int_{\partial B_R}K(x,y)u(y)d\sigma_y \qquad \text{for any } |x|<R.
\]
This is called the Poisson integral formula. Then we have (K5) easily by taking
$u\equiv 1$.
\end{quote}

The following result yields the existence of harmonic functions in balls with the
prescribed Dirichlet boundary value.

\begin{theorem}\label{hanlin:thm-poisson-dirichlet}\dcref{Theorem 1.24}\group{harmonic-functions}\source{han-lin-lecture-notes:page-0018}\uses{hanlin:cor-green-normal-derivative}
\emph{For $\varphi\in C(\partial B_R)$, the function $u$ defined by}
\begin{equation}
u(x)=\int_{\partial B_R}K(x,y)\varphi(y)d\sigma_y \qquad \text{for any } |x|<R
\tag{1}
\end{equation}
\emph{is smooth in $B_R$ and continuous up to $\partial B_R$ and satisfies}
\[
\left\{
\begin{aligned}
\Delta u &= 0 \qquad \text{in } \Omega\\
u &= \varphi \qquad \text{on } \partial \Omega.
\end{aligned}
\right.
\]
\end{theorem}

\begin{proof}
By (K1) and (K4), we conclude easily that $u$ defined by (1) is smooth
and harmonic in $B_R$. We only need to prove the continuity of $u$ up to the boundary
$\partial B_R$. Let $x_0\in \partial B_R$ and $x\in B_R$. By (K5), we have
\[
u(x)-\varphi(x_0)=\int_{|y|=R}K(x,y)(\varphi(y)-\varphi(x_0))d\sigma_y=I_1+I_2,
\]
where
\[
I_1=\int_{|y-x_0|<\delta,\,|y|=R}\cdots,\qquad
I_2=\int_{|y-x_0|>\delta,\,|y|=R}\cdots,
\]
for a positive constant $\delta$ to be determined. For any given $\varepsilon>0$, we choose $\delta=
\delta(\varepsilon)>0$ so small that
\[
|\varphi(y)-\varphi(x_0)|<\varepsilon \qquad for any |y-x_0|<\delta,\ |y|=R,
\]
by the continuity of $\varphi$. Then $|I_1|\le \varepsilon$ by (K2) and (K5). Let $M=\sup_{\partial B_R}|\varphi|$. By
(K3), we find a $\delta'$ such that
\[
K(x,y)\le \frac{\varepsilon}{2M\omega_n R^{n-1}} \qquad for any |x-x_0|<\delta',\ |y-x_0|>\delta,
\]
where $\delta'$ depends on $\varepsilon$ and $\delta=\delta(\varepsilon)$, and hence only on $\varepsilon$. Then $|I_2|<\varepsilon$. Hence
\[
|u(x)-\varphi(x_0)|<2\varepsilon \qquad for any |x-x_0|<\delta',\ |x|<R.
\]
This shows that $u$ is continuous at the boundary point $x_0$ and hence completes the proof.
\end{proof}

\begin{remark}\label{hanlin:remark-poisson-mean-value}\dcref{Remark 1.25}\group{harmonic-functions}\source{han-lin-lecture-notes:page-0019}\uses{hanlin:thm-poisson-dirichlet}
In the Poisson integral formula, by letting $x = 0$, we have
\[
u(0) = \frac{1}{\omega_n R^{n-1}} \int_{\partial B_R} u(y)\,d\sigma_y,
\]
which is the mean value property.
\end{remark}

\begin{lemma}\label{hanlin:lemma-harnack-ball}\dcref{Lemma 1.26}\group{harmonic-functions}\source{han-lin-lecture-notes:page-0019}\uses{hanlin:thm-poisson-dirichlet}
Suppose $u$ is harmonic in $B_R(x_0)$ and $u \ge 0$. Then
\[
\left(\frac{R}{R+r}\right)^{n-2}\frac{R-r}{R+r}u(x_0) \le u(x) \le \left(\frac{R}{R-r}\right)^{n-2}\frac{R+r}{R-r}u(x_0),
\]
where $r = |x-x_0| < R$.
\end{lemma}

\begin{proof}
We assume $x_0 = 0$ and $u \in C(\bar B_R)$. Note that $u$ is given by the Poisson integral formula
\[
u(x) = \frac{1}{\omega_n R}\int_{\partial B_R}\frac{R^2-|x|^2}{|y-x|^n}u(y)\,d\sigma_y.
\]
Since $R-|x| \le |y-x| \le R+|x|$ for $|y|=R$, we have
\[
\frac{1}{\omega_n R}\cdot\frac{R-|x|}{R+|x|}\left(\frac{1}{R+|x|}\right)^{n-2}\int_{\partial B_R}u(y)\,d\sigma_y \le u(x)
\]
\[
\le \frac{1}{\omega_n R}\cdot\frac{R+|x|}{R-|x|}\left(\frac{1}{R-|x|}\right)^{n-2}\int_{\partial B_R}u(y)\,d\sigma_y.
\]
The mean value property implies
\[
u(0) = \frac{1}{\omega_n R^{n-1}}\int_{\partial B_R}u(y)\,d\sigma_y.
\]
This finishes the proof.
\end{proof}

\begin{corollary}\label{hanlin:corollary-liouville-poisson}\dcref{Corollary 1.27}\group{harmonic-functions}\source{han-lin-lecture-notes:page-0019}\uses{hanlin:lemma-harnack-ball}
If $u$ is a harmonic function in $\mathbb{R}^n$ and bounded above or below, then $u \equiv \mathrm{const}$.
\end{corollary}

\begin{proof}
We assume $u \ge 0$ in $\mathbb{R}^n$. Take any point $x \in \mathbb{R}^n$ and apply Lemma 1.26 to any ball $B_R(0)$ with $R > |x|$. We obtain
\[
\left(\frac{R}{R+|x|}\right)^{n-2}\frac{R-|x|}{R+|x|}u(0) \le u(x) \le \left(\frac{R}{R-|x|}\right)^{n-2}\frac{R+|x|}{R-|x|}u(0).
\]
This implies $u(x)=u(0)$ by letting $R\to +\infty$.
\end{proof}

\begin{theorem}\label{hanlin:thm-removable-singularity}\dcref{Theorem 1.28}\group{harmonic-functions}\source{han-lin-lecture-notes:page-0019,han-lin-lecture-notes:page-0020}\uses{hanlin:thm-mean-value-max}
Suppose $u$ is harmonic in $B_R \setminus \{0\}$ and satisfies
\[
u(x)=
\begin{cases}
o(\log |x|), & n=2 \\
o(|x|^{2-n}), & n\ge 3
\end{cases}
\quad\text{as } |x|\to 0.
\]
Then $u$ can be defined at $0$ so that it is $C^2$ and harmonic in $B_R$.
\end{theorem}

\begin{proof}
Assume $u$ is continuous in $0 < |x| \le R$. Let v solve
\[
\left\{
\begin{aligned}
\Delta v &= 0 \text{ in } B_R \\
v &= u \text{ on } \partial B_R.
\end{aligned}
\right.
\]
We prove $u = v$ in $B_R \setminus \{0\}$. Set $w = v - u$ in $B_R \setminus \{0\}$ and $M_r = \max_{\partial B_r} |w|$. We only consider the case $n \ge 3$. It is obvious that
\[
|w(x)| \le M_r \cdot \frac{r^{n-2}}{|x|^{n-2}} \text{ on } \partial B_r.
\]
Note that $w$ and $\frac{1}{|x|^{n-2}}$ are harmonic in $B_R \setminus B_r$. Hence the maximum principle implies
\[
|w(x)| \le M_r \cdot \frac{r^{n-2}}{|x|^{n-2}} \text{ for any } x \in B_R \setminus B_r,
\]
where
\[
M_r = \max_{\partial B_r} |v - u| \le \max_{\partial B_r} |v| + \max_{\partial B_r} |u| \le M + \max_{\partial B_r} |u|,
\]
with $M = \max_{\partial B_R} |u|$. Then we have for each fixed $x \ne 0$
\[
|w(x)| \le \frac{r^{n-2}}{|x|^{n-2}} M + \frac{1}{|x|^{n-2}} r^{n-2} \max_{\partial B_r} |u| \to 0 \text{ as } r \to 0.
\]
This implies $w = 0$ in $B_R \setminus \{0\}$.
\end{proof}

\section*{1.3. Maximum Principles}

In this section, we use the maximum principle to derive the interior gradient estimate and the Harnack inequality. We first give another proof of the maximum principle without using mean value properties.

\begin{theorem}\label{hanlin:thm-subharmonic-max}\dcref{Theorem 1.29}\group{harmonic-functions}\source{han-lin-lecture-notes:page-0020}
Suppose $u \in C^2(B_1) \cap C(\bar B_1)$ is a subharmonic function in $B_1$, i.e., $\Delta u \ge 0$. Then
\[
\sup_{B_1} u \le \sup_{\partial B_1} u.
\]
\end{theorem}

\begin{proof}
For any $\varepsilon > 0$, we consider $u_\varepsilon(x) = u(x) + \varepsilon |x|^2$ in $B_1$. Then a simple calculation yields
\[
\Delta u_\varepsilon = \Delta u + 2n \varepsilon \ge 2n \varepsilon > 0.
\]
It is easy to see, by a contradiction argument, that $u_\varepsilon$ can not have an interior maximum. This implies in particular
\[
\sup_{B_1} u_\varepsilon \le \sup_{\partial B_1} u_\varepsilon.
\]
Therefore, we have
\[
\sup_{B_1} u \le \sup_{B_1} u_\varepsilon \le \sup_{\partial B_1} u + \varepsilon.
\]
We finish the proof by letting $\varepsilon \to 0$.
\end{proof}

The following result is called the comparison principle.

\begin{corollary}\label{hanlin:corollary-comparison-principle}\dcref{Corollary 1.30}\group{harmonic-functions}\source{han-lin-lecture-notes:page-0021}\uses{hanlin:thm-subharmonic-max}
Suppose $u,v \in C^2(B_1)\cap C(\overline{B}_1)$ satisfy
\[
\Delta u \ge \Delta v \quad \text{in } B_1
\]
\[
u \le v \quad \text{on } \partial B_1.
\]
Then $u \le v$ in $B_1$.
\end{corollary}

\begin{remark}\label{hanlin:remark-maxprin-domain}\dcref{Remark 1.31}\group{harmonic-functions}\source{han-lin-lecture-notes:page-0021}\uses{hanlin:thm-subharmonic-max,hanlin:corollary-comparison-principle}
Theorem 1.29 and Corollary 1.30 still hold if $B_1$ is replaced by
any bounded domain.
\end{remark}

\begin{theorem}\label{hanlin:thm-interior-gradient}\dcref{Theorem 1.32}\group{harmonic-functions}\source{han-lin-lecture-notes:page-0021}\uses{hanlin:thm-subharmonic-max}
Suppose $u$ is harmonic in $B_1$. Then
\[
\sup_{B_{1/2}} |Du| \le c\sup_{\partial B_1} |u|,
\]
where $c$ is a positive constant depending only on $n$.
\end{theorem}

By Theorem 1.32, we have for any $\alpha \in [0,1]$
\[
|u(x)-u(y)| \le c|x-y|^\alpha \sup_{\partial B_1} |u| \qquad \text{for any } x,y \in B_{1/2},
\]
where $c$ is a positive constant depending only on $n$.

\begin{proof}
A direct calculation shows
\[
\Delta(|Du|^2) = 2\sum_{i,j=1}^n (D_{ij}u)^2 + 2\sum_{i=1}^n D_i u D_i(\Delta u) = 2\sum_{i,j=1}^n (D_{ij}u)^2,
\]
where we used $\Delta u = 0$ in $B_1$. Hence $|Du|^2$ is a subharmonic function. To get
interior estimates we need a cut-off function. For any $\varphi \in C_0^1(B_1)$, we have
\[
\Delta(\varphi|Du|^2) = (\Delta\varphi)|Du|^2 + 4\sum_{i,j=1}^n D_i\varphi D_j u D_{ij}u + 2\varphi\sum_{i,j=1}^n (D_{ij}u)^2.
\]

By taking $\varphi=\eta^2$ for some $\eta\in C_0^1(B_1)$ with $\eta\equiv 1$ in $B_{1/2}$, we obtain by the Cauchy
inequality
\[
\Delta(\eta^2|Du|^2)
\]
\[
=2\eta \Delta\eta |Du|^2 + 2|D\eta|^2|Du|^2 + 8\eta \sum_{i,j=1}^n D_i\eta D_j u D_{ij}u + 2\eta^2 \sum_{i,j=1}^n (D_{ij}u)^2
\]
\[
\ge \bigl(2\eta\Delta\eta - 6|D\eta|^2\bigr)|Du|^2 \ge -C|Du|^2,
\]
where $C$ is a positive constant depending only on $\eta$. Note
\[
\Delta(u^2) = 2|Du|^2 + 2u\Delta u = 2|Du|^2,
\]
since $u$ is harmonic. By taking $\alpha$ large enough, we get
\[
\Delta(\eta^2|Du|^2 + \alpha u^2) \ge 0.
\]
We apply Theorem 1.29, the maximum principle, to get the result. $\square$
\end{proof}

Next we derive the Harnack inequality.

\begin{lemma}\label{hanlin:lemma-log-gradient}\dcref{Lemma 1.33}\group{harmonic-functions}\source{han-lin-lecture-notes:page-0022}
Suppose $u$ is a positive harmonic function in $B_1$. Then
\[
\sup_{B_{\frac{1}{2}}} |D \log u| \le c,
\]
where $c$ is a positive constant depending only on $n$.
\end{lemma}

\begin{proof}
Set $v = \log u$. Then a direct calculation shows
\[
\Delta v = -|Dv|^2.
\]

We need an interior gradient estimate on $v$. Set $w = |Dv|^2$. Then we get
\[
\Delta w + 2 \sum_{i=1}^n D_i v D_i w = 2 \sum_{i,j=1}^n (D_{ij}v)^2.
\]

As before, we need to introduce a cut-off function. First note
\begin{equation}
\sum_{i,j=1}^n (D_{ij}v)^2 \ge \sum_{i=1}^n (D_{ii}v)^2 \ge \frac{1}{n}(\Delta v)^2 = \frac{|Dv|^4}{n} = \frac{w^2}{n}.
\tag{1}
\end{equation}

Take a nonnegative function $\varphi \in C_0^1(B_1)$. We obtain by the Cauchy inequality
\begin{align*}
\Delta(\varphi w) + 2 \sum_{i=1}^n D_i v D_i(\varphi w)
&= 2\varphi \sum_{i,j=1}^n (D_{ij}v)^2 + 4 \sum_{i,j=1}^n D_i\varphi D_j v D_{ij}v + 2w \sum_{i=1}^n D_i\varphi D_i v + (\Delta\varphi)w \\
&\ge \varphi \sum_{i,j=1}^n (D_{ij}v)^2 - 2|D\varphi||Dv|^3 - \left(|\Delta\varphi| + C\frac{|D\varphi|^2}{\varphi}\right)|Dv|^2,
\end{align*}
if we choose $\varphi$ such that $|D\varphi|^2/\varphi$ is bounded in $B_1$. Choose $\varphi = \eta^4$ for some $\eta \in C_0^1(B_1)$. For such a fixed $\eta$, we obtain by (1)
\begin{align*}
\Delta(\eta^4 w) + 2 \sum_{i=1}^n D_i v D_i(\eta^4 w)
&\ge \frac{1}{n}\eta^4 |Dv|^4 - Cn^3|D\eta||Dv|^3 - 4n^2(n\Delta\eta + C|D\eta|^2)|Dv|^2 \\
&\ge \frac{1}{n}\eta^4 |Dv|^4 - Cn^3|Dv|^3 - C\eta^2|Dv|^2,
\end{align*}
where $C$ is a positive constant depending only on $n$ and $\eta$. Hence we get by the Cauchy inequality
\[
\Delta(\eta^4 w) + 2 \sum_{i=1}^n D_i v D_i(\eta^4 w) \ge \frac{1}{n}\eta^4 w^2 - C,
\]
where $C$ is a positive constant depending only on $n$ and $\eta$.

Suppose $\eta^4 w$ attains its maximum at $x_0 \in B_1$. Then $D(\eta^4 w)=0$ and
\[
\Delta(\eta^4 w) \le 0 \quad \text{at } x_0.
\]
Hence
\[
\eta^4 w^2(x_0) \le C(n,\eta).
\]
If $w(x_0) \ge 1$, then $\eta^4 w(x_0) \le C(n)$. Otherwise $\eta^4 w(x_0) \le w(x_0) \le \eta^4(x_0)$. In both cases we conclude
\[
\eta^4 w \le C(n,\eta) \quad \text{in } B_1.
\]
\end{proof}

\begin{corollary}\label{hanlin:corollary-harnack-ball}\dcref{Corollary 1.34}\group{harmonic-functions}\source{han-lin-lecture-notes:page-0023}\uses{hanlin:lemma-log-gradient}
Suppose $u$ is a nonnegative harmonic function in $B_1$. Then
\[
u(x_1) \le c u(x_2) \text{ for any } x_1, x_2 \in B_{\frac{1}{2}},
\]
where $c$ is a positive constant depending only on $n$.
\end{corollary}

\begin{proof}
We may assume $u > 0$ in $B_1$. For any $x_1, x_2 \in B_{\frac{1}{2}}$, by a simple integration we obtain with Lemma 1.33
\[
\log \frac{u(x_1)}{u(x_2)} \le |x_1 - x_2| \int_0^1 |D \log u(tx_2 + (1 - t)x_1)|\,dt \le C|x_1 - x_2|.
\]
This finishes the proof.
\end{proof}

\begin{theorem}\label{hanlin:thm-hopf}\dcref{Theorem 1.35}\group{harmonic-functions}\source{han-lin-lecture-notes:page-0023,han-lin-lecture-notes:page-0024}\uses{hanlin:thm-subharmonic-max,hanlin:corollary-harnack-ball}
Suppose $u \in C(\bar{B}_1)$ is a harmonic function in $B_1$. If $u(x) < u(x_0)$ for any $x \in \bar{B}_1$ and some $x_0 \in \partial B_1$, then
\[
\frac{\partial u}{\partial n}(x_0) \ge c(u(x_0) - u(0)),
\]
where $c$ is a positive constant depending only on $n$.
\end{theorem}

\begin{proof}
Consider a positive function $v$ in $B_1$ defined by
\[
v(x) = e^{-\alpha |x|^2} - e^{-\alpha}.
\]
It is easy to see
\[
\Delta v(x) = e^{-\alpha |x|^2}(-2\alpha n + 4\alpha^2|x|^2) > 0 \text{ for any } |x| \ge \frac{1}{2},
\]
if $\alpha \ge 2n+1$. Hence for such a fixed $\alpha$, the function $v$ is subharmonic in the region
\[
A = B_1 \setminus B_{\frac{1}{2}}.
\]
Now define for $\varepsilon > 0$
\[
h_\varepsilon(x) = u(x) - u(x_0) + \varepsilon v(x).
\]
This is also a subharmonic function, i.e., $\Delta h_\varepsilon > 0$ in $A$. Obviously $h_\varepsilon < 0$ on $\partial B_1$
and $h_\varepsilon(x_0) = 0$. Since $u(x) < u(x_0)$ for $|x| = \frac{1}{2}$, we take $\varepsilon > 0$ small such that
$h_\varepsilon(x) < 0$ for $|x| = \frac{1}{2}$. Therefore by Theorem 1.29, $h_\varepsilon$ assumes at the point $x_0$
its maximum in $A$. This implies
\[
\frac{\partial h_\varepsilon}{\partial n}(x_0) \ge 0,
\]
or
\[
\frac{\partial u}{\partial n}(x_0) \ge -\varepsilon \frac{\partial v}{\partial n}(x_0) = 2\alpha \varepsilon e^{-\alpha} > 0.
\]
Note that we only used the subharmonicity of $u$ so far. We estimate $\varepsilon$ as follows.
Set $w(x) = u(x_0) - u(x) > 0$ in $B_1$. Obviously $w$ is a harmonic function in $B_1$. By
Corollary 1.34, the Harnack inequality, we have
\[
\inf_{B_{\frac{1}{2}}} w \ge c w(0),
\]
or
\[
\max_{B_{\frac{1}{2}}} u \le u(x_0) - c(u(x_0) - u(0)),
\]
where $c$ is a positive constant depending only on $n$. Hence we take
\[
\varepsilon = \delta c\bigl(u(x_0)-u(0)\bigr),
\]
for $\delta$ small, depending on $n$. This finishes the proof.
\hfill$\square$
\end{proof}

To finish this section, we prove a result on the global H\"older continuity for
harmonic functions.

\begin{lemma}\label{hanlin:lemma-holder-boundary}\dcref{Lemma 1.36}\group{harmonic-functions}\source{han-lin-lecture-notes:page-0024,han-lin-lecture-notes:page-0025}\uses{hanlin:corollary-comparison-principle,hanlin:thm-interior-gradient}
Suppose $u \in C(\overline{B_1})$ is a harmonic function in $B_1$ with $u=\varphi$ on
$\partial B_1$. If $\varphi \in C^\alpha(\partial B_1)$ for some $\alpha \in (0,1)$, then
$u \in C^{\alpha/2}(\overline{B_1})$ and
\[
\|u\|_{C^{\alpha/2}(\overline{B_1})} \le c\|\varphi\|_{C^\alpha(\partial B_1)},
\]
where $c$ is a positive constant depending only on $n$ and $\alpha$.
\end{lemma}

\begin{proof}
First, the maximum principle implies that
\[
\inf_{\partial B_1}\varphi \le u \le \sup_{\partial B_1}\varphi \qquad \text{in } B_1.
\]
Next, we claim for any $x_0 \in \partial B_1$
\begin{equation}\tag{1}
\sup_{x \in B_1}\frac{|u(x)-u(x_0)|}{|x-x_0|^{\alpha/2}}
\le 2^{\alpha/2}\sup_{x \in \partial B_1}\frac{|\varphi(x)-\varphi(x_0)|}{|x-x_0|^\alpha}.
\end{equation}
Lemma 1.36 follows easily from (1).

For any $x,y \in B_1$, set $d_x = \operatorname{dist}(x,\partial B_1)$ and
$d_y = \operatorname{dist}(y,\partial B_1)$, and assume $d_y \le d_x$. Take
$x_0,y_0 \in \partial B_1$ such that $|x-x_0| = d_x$ and $|y-y_0| = d_y$. Assume first
that $|x-y| \le d_x/2$. Then $y \in \overline{B}_{d_x/2}(x) \subset B_{d_x}(x) \subset B_1$.
We apply Theorem 1.32 (scaled version) to $u-u(x_0)$ in $B_{d_x}(x)$ and get by (1)
\[
\frac{d_x^{\alpha/2}|u(x)-u(y)|}{|x-y|^{\alpha/2}}
\le C\bigl|u-u(x_0)\bigr|_{L^\infty(B_{d_x}(x))}
\le C d_x^{\alpha/2}\|\varphi\|_{C^\alpha(\partial B_1)}.
\]
Hence, we obtain
\[
|u(x)-u(y)| \le C|x-y|^{\alpha/2}\|\varphi\|_{C^\alpha(\partial B_1)}.
\]

Assume now that $d_y \le d_x \le 2|x-y|$. Then by (1) again we have
\[
|u(x)-u(y)| \le |u(x)-u(x_0)| + |u(x_0)-u(y_0)| + |u(y_0)-u(y)|
\]
\[
\le C\bigl(d_x^{\alpha/2} + |x_0-y_0|^{\alpha/2} + d_y^{\alpha/2}\bigr)\|\varphi\|_{C^\alpha(\partial B_1)}
\]
\[
\le C|x-y|^{\alpha/2}\|\varphi\|_{C^\alpha(\partial B_1)},
\]
since $|x_0-y_0| \le d_x + |x-y| + d_y \le 5|x-y|$.

To prove (1), we assume $B_1 = B_1((1,0,\dots,0))$, $x_0=0$ and $\varphi(0)=0$. Define
$K=\sup_{x\in \partial B_1}|\varphi(x)|/|x|^\alpha$. Note $|x|^2=2x_1$ for $x\in \partial B_1$.
Then
\[
\varphi(x) \le K|x|^\alpha \le 2^{\alpha/2}Kx_1^{\alpha/2}\qquad \text{for any }x\in\partial B_1.
\]
Define $v(x)=2^{\alpha/2}Kx_1^{\alpha/2}$ in $B_1$. Then we have
\[
\Delta v(x)=2^{\alpha/2}K\cdot \frac{\alpha}{2}\left(\frac{\alpha}{2}-1\right)x_1^{\alpha/2-2}<0
\qquad \text{in }B_1.
\]
Corollary 1.30 implies
\[
u(x)\le v(x)=2^{\alpha/2}Kx_1^{\alpha/2}\le 2^{\alpha/2}K|x|^{\alpha/2}
\qquad \text{for any }x\in B_1.
\]

Considering $-u$ similarly, we get
\[
|u(x)| \le 2^{\alpha/2}K|x|^{\alpha/2}\qquad \text{for any }x\in B_1.
\]
This proves (1).
\hfill$\square$
\end{proof}

\subsection*{1.4. Energy Methods}

In this section, we discuss harmonic functions by using energy methods. In
general, we assume throughout this section that $a_{ij}\in C(B_1)$ satisfies
\[
\lambda |\xi|^2 \le a_{ij}(x)\xi_i\xi_j \le \Lambda |\xi|^2\qquad \text{for any }x\in B_1\text{ and }\xi\in\mathbb{R}^n,
\]
for some positive constants $\lambda$ and $\Lambda$. We consider the function $u\in C^1(B_1)$ satisfying
\[
\int_{B_1} a_{ij}\D_i u\D_j \varphi = 0\qquad \text{for any }\varphi\in C_0^1(B_1).
\]
It is easy to check by integration by parts that harmonic functions satisfy the above
equation for $a_{ij}=\delta_{ij}$.

The following result is referred to as the Caccioppoli inequality.

\begin{lemma}\label{hanlin:lemma-caccioppoli}\dcref{Lemma 1.37}\group{harmonic-functions}\source{han-lin-lecture-notes:page-0025}
\emph{Suppose $u\in C^1(B_1)$ satisfies}
\[
\int_{B_1} a_{ij}\D_i u\D_j \varphi = 0\qquad \text{for any }\varphi\in C_0^1(B_1).
\]
\emph{Then for any function $\eta\in C_0^1(B_1)$}
\[
\int_{B_1} \eta^2|\D u|^2 \le c\int_{B_1} |\D\eta|^2 u^2,
\]
\emph{where $c$ is a positive constant depending only on $\lambda$ and $\Lambda$.}
\end{lemma}

\begin{proof}
For any $\eta\in C_0^1(B_1)$, set $\varphi=\eta^2u$. Then we have
\[
\lambda \int_{B_1} \eta^2|\D u|^2 \le \Lambda \int_{B_1} \eta|u||\D\eta||\D u|.
\]
A simple application of the Cauchy inequality yields the result.
\hfill$\square$
\end{proof}

\begin{corollary}\label{hanlin:cor-caccioppoli-balls}\dcref{Corollary 1.38}\group{harmonic-functions}\source{han-lin-lecture-notes:page-0025}\uses{hanlin:lemma-caccioppoli}
\emph{Let $u$ be as in Lemma 1.37. Then for any $0\le r < R\le 1$}
\[
\int_{B_r} |\D u|^2 \le \frac{C}{(R-r)^2}\int_{B_R} u^2,
\]
\emph{where $C$ is a positive constant depending only on $\lambda$ and $\Lambda$.}
\end{corollary}

\begin{proof}
Take $\eta$ such that $\eta=1$ on $B_r$, $\eta=0$ outside $B_R$ and $|\D\eta|\le 2(R-r)^{-1}$.
\hfill$\square$
\end{proof}

\begin{corollary}\label{hanlin:cor-energy-decay}\dcref{Corollary 1.39}\group{harmonic-functions}\source{han-lin-lecture-notes:page-0025,han-lin-lecture-notes:page-0026}\uses{hanlin:lemma-caccioppoli}
\emph{Let u be as in Lemma 1.37. Then for any $0<R\le 1$}
\[
\int_{B_{R/2}} u^2 \le \theta \int_{B_R} u^2,
\]
\emph{and}
\[
\int_{B_{R/2}} |\D u|^2 \le \theta \int_{B_R} |\D u|^2,
\]
\emph{where $\theta\in(0,1)$ is a positive constant depending only on $n,\lambda$ and $\Lambda$.}
\end{corollary}

\begin{proof}
Take an $\eta \in C_0^1(B_R)$ with $\eta = 1$ on $B_{R/2}$ and $|D\eta| \le 2R^{-1}$. Then Lemma 1.37 yields
\[
\int_{B_R} |D(\eta u)|^2 \le C \int_{B_R} |D\eta|^2 u^2 \le \frac{C}{R^2} \int_{B_R \setminus B_{R/2}} u^2,
\]
since $D\eta = 0$ in $B_{R/2}$. Hence, by the Poincar\'e inequality, we get
\[
\int_{B_{R/2}} (\eta u)^2 \le C R^2 \int_{B_R} |D(\eta u)|^2.
\]

Therefore we obtain
\[
\int_{B_{R/2}} u^2 \le C \int_{B_R \setminus B_{R/2}} u^2,
\]
and hence
\[
(C + 1) \int_{B_{R/2}} u^2 \le C \int_{B_R} u^2.
\]

For the second inequality, observe that Lemma 1.37 holds for $u - a$ for an arbitrary constant $a$. Then as before, we have
\[
\int_{B_R} \eta^2 |Du|^2 \le C \int_{B_R} |D\eta|^2 (u - a)^2 \le \frac{C}{R^2} \int_{B_R \setminus B_{R/2}} (u - a)^2.
\]
The Poincar\'e inequality implies with $a = |B_R \setminus B_{R/2}|^{-1} \int_{B_R \setminus B_{R/2}} u$
\[
\int_{B_R \setminus B_{R/2}} (u - a)^2 \le c(n) R^2 \int_{B_R \setminus B_{R/2}} |Du|^2.
\]
Hence we obtain
\[
\int_{B_{R/2}} |Du|^2 \le C \int_{B_R \setminus B_{R/2}} |Du|^2,
\]
and hence
\[
(C + 1) \int_{B_{R/2}} |Du|^2 \le C \int_{B_R} |Du|^2.
\]

This finishes the proof.
\end{proof}

\begin{remark}\label{hanlin:remark-finite-energy-liouville}\dcref{Remark 1.40}\group{harmonic-functions}\source{han-lin-lecture-notes:page-0026}\uses{hanlin:cor-energy-decay}
Corollary 1.39 implies, in particular, that a harmonic function in $\mathbb{R}^n$ with a finite $L^2$-norm is identically zero and that a harmonic function in $\mathbb{R}^n$ with a finite Dirichlet integral is constant.
\end{remark}

\begin{remark}\label{hanlin:remark-iterated-energy-decay}\dcref{Remark 1.41}\group{harmonic-functions}\source{han-lin-lecture-notes:page-0026}\uses{hanlin:lemma-caccioppoli,hanlin:cor-energy-decay}
By iterating the result in Corollary 1.39, we have the following estimates. Let $u$ be as in Lemma 1.37. Then for any $0 < \rho < r \le 1$
\[
\int_{B_\rho} u^2 \le C \left(\frac{\rho}{r}\right)^\mu \int_{B_r} u^2,
\]
and
\[
\int_{B_\rho} |Du|^2 \le C \left(\frac{\rho}{r}\right)^\mu \int_{B_r} |Du|^2,
\]
for some positive constant $\mu$ depending only on $n$, $\lambda$ and $A$. Later on, we will prove that we can take $\mu \in (n - 2,n)$ in the second inequality. For harmonic functions, we have better results.
\end{remark}

\begin{lemma}\label{hanlin:lemma-constant-coeff-compactness}\dcref{Lemma 1.42}\group{harmonic-functions}\source{han-lin-lecture-notes:page-0027,han-lin-lecture-notes:page-0028}\uses{hanlin:cor-caccioppoli-balls}
Suppose $\{a_{ij}\}$ is a constant positive definite matrix with
\[
\lambda |\xi|^2 \le a_{ij}\xi_i\xi_j \le \Lambda |\xi|^2 \qquad \text{for any } \xi \in \mathbb{R}^n,
\]
for some constants $0 < \lambda \le \Lambda$. Suppose $u \in C^1(B_1)$ satisfies
\begin{equation}
\int_{B_1} a_{ij}D_i u D_j \varphi = 0 \qquad \text{for any } \varphi \in C_0^1(B_1).
\end{equation}
Then for any $0 < \rho \le r$
\begin{equation}
\int_{B_\rho} |u|^2 \le c \left(\frac{\rho}{r}\right)^n \int_{B_r} |u|^2,
\end{equation}
and
\begin{equation}
\int_{B_\rho} |u-u_\rho|^2 \le c \left(\frac{\rho}{r}\right)^{n+2} \int_{B_r} |u-\bar{u}_r|^2,
\end{equation}
where $c$ is a positive constant depending only on $\lambda$ and $\Lambda$ and $\bar{u}_r$ denotes the average
of $u$ in $B_r$.
\end{lemma}

\begin{proof}
By a simple dilation, we need only consider $r = 1$. We only prove (2)
and (3) for $\rho \in (0,\frac12]$, since they are trivial for $\rho \in (\frac12,1]$.
Now we claim
\[
\lVert u \rVert_{L^\infty(B_{1/2})}^2 + \lVert Du \rVert_{L^\infty(B_{1/2})}^2 \le c \int_{B_1} |u|^2,
\]
where $c$ is a positive constant depending only on $\lambda$ and $\Lambda$. This implies for $\rho \in (0,\frac12]$
\[
\int_{B_\rho} |u|^2 \le \rho^n \lVert u \rVert_{L^\infty(B_{1/2})}^2 \le c\rho^n \int_{B_1} |u|^2,
\]
and
\[
\int_{B_\rho} |u-\bar{u}_\rho|^2 \le \int_{B_\rho} |u-u(0)|^2 \le \rho^{n+2}\lVert Du \rVert_{L^\infty(B_{1/2})}^2 \le c\rho^{n+2}\int_{B_1} |u|^2.
\]
If $u$ is a solution of (1), so is $u-u_1$. With $u$ replaced by $u-u_1$ in the above
inequality, we have
\[
\int_{B_\rho} |u-\bar{u}_\rho|^2 \le c\rho^{n+2}\int_{B_1} |u-u_1|^2.
\]

We now present two methods to prove the claim.

Method 1. By a rotation, we may assume $(a_{ij})$ is a diagonal matrix. Hence
(1) becomes
\[
\sum_{i=1}^n \lambda_i D_{ii}u = 0,
\]
with $0 < \lambda \le \lambda_i \le \Lambda$ for $i = 1,\cdots,n$. It is easy to see that there exists a constant
$r_0 \in (0,\frac12)$, depending only on $\lambda$ and $\Lambda$, such that for any $x_0 \in B_{1/2}$ the rectangle
\[
\left\{x; \frac{|x_i-x_{0i}|}{\sqrt{\lambda_i}} < r_0 \right\}
\]
is contained in $B_1$. Consider the change the coordinate
\[
x_i \mapsto y_i = \frac{x_i}{\sqrt{\lambda_i}},
\]
and set
\[
v(y)=u(x).
\]
Then $v$ is harmonic in $\{y;\sum_{i=1}^n \lambda_i y_i^2 < 1\}$. In the ball
\[
\{y;|y-y_0|<r_0\},
\]
we use the interior estimates to yield
\[
|v(y_0)|^2+|Dv(y_0)|^2 \le c \int_{B_{r_0}(y_0)} v^2 \le c \int_{\{\sum_{i=1}^n \lambda_i y_i^2<1\}} v^2,
\]
where $c$ is a positive constant depending only on $\lambda$ and $\Lambda$. We transform back to $u$
to get
\[
|u(x_0)|^2+|Du(x_0)|^2 \le c \int_{|x|<1} u^2.
\]

\textit{Method 2.} If $u$ is a solution of (1), so are any derivatives of $u$. By applying
Corollary 1.38 to derivatives of $u$, we conclude for any positive integer $k$
\[
\|u\|_{H^k(B_{1/2})} \le c\|u\|_{L^2(B_1)},
\]
where $c$ is a positive constant depending only on $n,\lambda$ and $\Lambda$. By taking
$k=\left[\frac{n}{2}\right]+1$, we have $H^k(B_{1/2})$ continuously embedded in $C^1(\overline{B_{1/2}})$. We then get
\[
\|u\|_{L^\infty(B_{1/2})}+\|Du\|_{L^\infty(B_{1/2})} \le c\|u\|_{L^2(B_1)}.
\]
This finishes the proof.
\end{proof}
